\documentclass[11pt,a4paper]{article}

\usepackage{fontspec}
\usepackage[margin=20mm]{geometry}
\usepackage{amsmath}
\usepackage{booktabs,tabularx,array}
\usepackage{xcolor}
\usepackage{hyperref}
\usepackage{enumitem}

\defaultfontfeatures{Ligatures=TeX,Scale=MatchLowercase}
\IfFontExistsTF{Noto Sans CJK KR}{
  \setmainfont{Noto Sans CJK KR}[
    ItalicFeatures={FakeSlant=0.18},
    BoldItalicFeatures={FakeSlant=0.18}
  ]
  \setsansfont{Noto Sans CJK KR}[
    ItalicFeatures={FakeSlant=0.18},
    BoldItalicFeatures={FakeSlant=0.18}
  ]
}{
  \IfFontExistsTF{Apple SD Gothic Neo}{
    \setmainfont{Apple SD Gothic Neo}[
      ItalicFeatures={FakeSlant=0.18},
      BoldItalicFeatures={FakeSlant=0.18}
    ]
    \setsansfont{Apple SD Gothic Neo}[
      ItalicFeatures={FakeSlant=0.18},
      BoldItalicFeatures={FakeSlant=0.18}
    ]
  }{
    \setmainfont{NanumGothic}[
      ItalicFeatures={FakeSlant=0.18},
      BoldItalicFeatures={FakeSlant=0.18}
    ]
    \setsansfont{NanumGothic}[
      ItalicFeatures={FakeSlant=0.18},
      BoldItalicFeatures={FakeSlant=0.18}
    ]
  }
}
\IfFontExistsTF{DejaVu Sans Mono}{
  \setmonofont{DejaVu Sans Mono}
}{
  \IfFontExistsTF{Menlo}{
    \setmonofont{Menlo}
  }{
    \setmonofont{Liberation Mono}
  }
}

\XeTeXlinebreaklocale "ko"
\XeTeXlinebreakskip = 0pt plus 1pt

\hypersetup{
  colorlinks=true,
  linkcolor=blue!50!black,
  urlcolor=blue!50!black,
  pdftitle={UUV MuJoCo v2.2 ellipsoid fluid parameter note},
  pdfauthor={OpenAI Codex}
}

\setlength{\parskip}{0.45em}
\setlength{\parindent}{0pt}
\renewcommand{\arraystretch}{1.18}
\newcolumntype{Y}{>{\raggedright\arraybackslash}X}
\newcommand{\code}[1]{\nolinkurl{#1}}

\title{UUV MuJoCo v2.2\\Ellipsoid Fluid Model 파라미터 정리본}
\author{}
\date{2026-04-07}

\begin{document}
\maketitle

\section{문서 목적}

이 문서는 \code{tank\_current\_scene.xml} 에서 사용하는 MuJoCo \textbf{ellipsoid fluid model} 이 어떤 힘을 계산하는지, 그리고 \textbf{5개의 \code{fluidcoef}} 가 각각 어떤 의미를 가지는지 정리한 설명본이다. 또한 current scene에서 \textbf{왜 center enclosure, lower body, plate proxy에 현재 값을 넣었는지}를 함께 기록한다.

기준 소스는 다음 두 가지다.
\begin{itemize}[leftmargin=2em]
  \item MuJoCo 공식 문서 \textit{Fluid forces}: \url{https://mujoco.readthedocs.io/en/3.1.1/computation/fluid.html}
  \item \code{uuv\_mujoco/v2.2/scenes/tank\_current\_scene.xml}
\end{itemize}

\section{MuJoCo ellipsoid model이 계산하는 힘}

MuJoCo의 ellipsoid fluid model은 geom 하나를 등가 타원체로 보고, 그 주변 유체가 만드는 힘과 토크를 아래와 같이 합산한다.

\begin{align*}
\mathbf{f}_{\text{ellipsoid}} &= \mathbf{f}_A + \mathbf{f}_D + \mathbf{f}_M + \mathbf{f}_K + \mathbf{f}_V \\
\mathbf{g}_{\text{ellipsoid}} &= \mathbf{g}_A + \mathbf{g}_D + \mathbf{g}_V
\end{align*}

여기서 각 항의 의미는 다음과 같다.
\begin{itemize}[leftmargin=2em]
  \item \(\mathbf{f}_A, \mathbf{g}_A\): added mass / added inertia
  \item \(\mathbf{f}_D, \mathbf{g}_D\): drag / angular drag
  \item \(\mathbf{f}_M\): Magnus lift
  \item \(\mathbf{f}_K\): Kutta lift
  \item \(\mathbf{f}_V, \mathbf{g}_V\): viscous resistance
\end{itemize}

\textbf{중요한 점은}, \code{fluidcoef} 의 5개 숫자는 이 전체 모델 중 \textbf{drag / angular drag / lift 계열을 조정하는 계수}라는 것이다. Added mass와 viscous resistance는 모델 형상, 유체 밀도 \(\rho\), 점성 \(\beta\) 에 의해 함께 계산되지만, \code{fluidcoef} 5개로 직접 스케일링하는 구조는 아니다.

\section{5개 파라미터의 의미와 수식}

MuJoCo 문서 기준으로 \code{fluidcoef = [0,1,2,3,4]} 는 아래 의미를 가진다.

\begin{center}
\begin{tabularx}{\linewidth}{@{}c l l Y@{}}
\toprule
인덱스 & 이름 & 기호 & 의미 \\
\midrule
0 & Blunt drag & \(C_{D,\text{blunt}}\) & 정면으로 부딪히는 면적에 대한 저항 \\
1 & Slender drag & \(C_{D,\text{slender}}\) & 길쭉한 형상 방향의 저항 보정 \\
2 & Angular drag & \(C_{D,\text{angular}}\) & 회전 시 저항 토크 \\
3 & Kutta lift & \(C_K\) & 받음각이 있는 얇은 형상에서 생기는 lift \\
4 & Magnus lift & \(C_M\) & 회전하며 이동하는 물체에서 생기는 lift \\
\bottomrule
\end{tabularx}
\end{center}

\subsection{0번: Blunt drag coefficient}

MuJoCo 문서의 drag force는 다음 형태로 정리된다.
\begin{align*}
\mathbf{f}_D
=
- \rho
\Big[
C_{D,\text{blunt}} A^{\text{proj}}_{\mathbf{v}}
+
C_{D,\text{slender}}(A_{\max} - A^{\text{proj}}_{\mathbf{v}})
\Big]
\|\mathbf{v}\| \mathbf{v}
\end{align*}

\(A^{\text{proj}}_{\mathbf{v}}\) 는 속도 방향에 대해 보이는 투영면적이다. 따라서 \textbf{0번 계수는 정면으로 물을 밀고 나갈 때의 압력성 저항}을 키우는 역할을 한다.

\subsection{1번: Slender drag coefficient}

1번은 위 식의 두 번째 항에 들어가며, \textbf{정면 투영면적만으로 설명되지 않는 길쭉한 형상의 drag} 를 보정한다. 길고 얇은 body는 blunt drag만으로는 맞추기 어렵기 때문에 이 계수가 같이 들어간다.

\subsection{2번: Angular drag coefficient}

회전 토크는 문서상 다음 형태로 쓴다.
\begin{align*}
\mathbf{g}_D
=
- \rho ~ \boldsymbol{\omega}
\Big(
\big[
C_{D,\text{angular}} \mathbf{I}_D
+
C_{D,\text{slender}}(\mathbf{I}_{\max}-\mathbf{I}_D)
\big]
\cdot \boldsymbol{\omega}
\Big)
\end{align*}

즉 2번은 \textbf{yaw, roll, pitch 회전 시 받는 저항 토크}를 키우는 항이다. current 모델에서 yaw 응답이 너무 가볍게 돌거나, 중앙 body가 회전에 비해 너무 가볍게 느껴질 때 가장 직접적으로 보는 값이다.

\subsection{3번: Kutta lift coefficient}

Kutta lift는 문서상 다음과 같이 쓴다.
\begin{align*}
\mathbf{f}_K
=
\frac{C_K \rho A^{\text{proj}}_{\mathbf{v}}}{\|\mathbf{v}\|}
\big(\mathbf{v} \times \mathbf{v}_\parallel\big)\times\mathbf{v}
\end{align*}

이는 \textbf{날개나 얇은 trailing edge가 있는 형상에서 받음각 때문에 발생하는 lift}에 해당한다. 일반적인 호버링 ROV에서는 의도하지 않은 측면력이나 상하 외란으로 이어질 수 있으므로 보수적으로 보는 것이 맞다.

\subsection{4번: Magnus lift coefficient}

Magnus lift는 문서상 다음과 같이 쓴다.
\begin{align*}
\mathbf{f}_M = C_M \rho V (\boldsymbol{\omega} \times \mathbf{v})
\end{align*}

즉 \textbf{회전하는 물체가 동시에 이동할 때 생기는 lift} 이다. 회전 실린더, 회전구, 강한 스핀 바디에서는 의미가 있지만, current UUV의 fluid proxy에는 기본적으로 의도한 효과가 아니다.

\section{왜 current scene에서는 이 값을 넣었는가}

현재 scene의 fluid proxy 값은 다음과 같다.

\begin{center}
\small
\begin{tabularx}{\linewidth}{@{}l c c c c c Y@{}}
\toprule
geom 그룹 & type & 0 & 1 & 2 & 3/4 & 설정 의도 \\
\midrule
Center enclosure & ellipsoid & 0.620 & 0.175 & 1.800 & 0 / 0 & 가장 큰 부피와 중앙 회전 저항을 담당 \\
Port / starboard lower body & ellipsoid & 0.520 & 0.150 & 1.300 & 0 / 0 & 전진 drag와 yaw 시 좌우 레버암 효과를 담당 \\
Top / bottom plates & box + ellipsoid fluid & 0.240 & 0.095 & 0.700 & 0 / 0 & 평평한 상하 면적 보강, heave 저항 보완 \\
\bottomrule
\end{tabularx}
\end{center}

\subsection{Center enclosure}

center enclosure의 current 값은 \code{0.620 0.175 1.800 0 0} 이다. 중앙의 가장 큰 부피 프록시이므로 \textbf{blunt drag를 가장 크게} 두었다. 또한 기체 중앙의 큰 body가 회전 시 가장 큰 물 저항을 만들어야 하므로 \textbf{angular drag도 가장 크게} 잡았다.

3번과 4번을 0으로 둔 이유는, 이 body에서 lift를 적극적으로 만들고 싶은 것이 아니라 \textbf{부피성 drag와 회전 감쇠를 표현하는 것}이 목적이기 때문이다.

\subsection{Port / starboard lower body}

좌우 lower body의 current 값은 \code{0.520 0.150 1.300 0 0} 이다. center보다 길쭉하고 단면이 작기 때문에 blunt, slender, angular 값을 조금 낮췄다. 그러나 이 두 body는 중심축에서 좌우로 떨어져 있으므로, \textbf{yaw 시에는 위치 자체가 레버암으로 작동해 전체 회전 감쇠에 크게 기여}한다.

즉 lower body는 자체 계수보다도 \textbf{좌우로 분리된 배치} 때문에 yaw 응답을 현실적으로 만드는 역할이 크다.

\subsection{Top / bottom plate}

top / bottom plate의 current 값은 \code{0.240 0.095 0.700 0 0} 이다. plate proxy를 추가한 이유는 lower body의 실제 상판 / 하판이 ellipsoid보다 더 평평하기 때문이다. 기존 ellipsoid-only 구성은 상하가 너무 둥글어서, \textbf{heave 방향 투영면적이 실제보다 작게 잡히고 z축 저항이 약하게 나올 가능성}이 있었다.

그래서 얇은 box proxy를 위아래에 추가해 평평한 상하 면적을 보강했다. 다만 이 보조 proxy가 전체 drag를 과하게 키우지 않도록 \textbf{center나 lower body보다 훨씬 작은 계수}에서 시작했다. 목적은 lift를 만드는 것이 아니라 \textbf{heave 감쇠를 보강하는 것}이다.

\section{왜 3번과 4번은 0으로 두었는가}

current UUV의 목적은 \textbf{전진 응답, yaw 감쇠, heave 감쇠를 실제와 비슷하게 맞추는 것}이다. 반면 Kutta lift와 Magnus lift는 주로 다음 상황에서 의미가 크다.

\begin{itemize}[leftmargin=2em]
  \item 얇은 날개, trailing edge, 받음각이 큰 비행체
  \item 강한 회전과 병진이 동시에 있는 스핀 바디
\end{itemize}

현재 UUV는 이런 lift를 의도적으로 활용하는 모델이 아니므로, 3번과 4번을 키우면 \textbf{원하지 않는 측면력이나 상하 외란}이 들어갈 가능성이 있다. 그래서 current scene은 \textbf{0,1,2번 drag 계수만 조정하고 3,4번은 0으로 유지}했다.

\section{한 줄 요약}

\begin{itemize}[leftmargin=2em]
  \item MuJoCo ellipsoid fluid model은 added mass, drag, angular drag, lift, viscous resistance를 함께 계산한다.
  \item 그러나 \code{fluidcoef} 5개는 그중 \textbf{drag / angular drag / lift 계열}을 조정하는 파라미터다.
  \item current scene에서는 \textbf{0,1,2번으로 전진 / yaw / heave 응답을 맞추고, 3,4번은 불필요한 lift를 막기 위해 0으로 둔다.}
  \item 새로 추가한 top / bottom plate proxy는 \textbf{평평한 상하 면적을 반영해 z축 유체저항을 보강}하기 위한 것이다.
\end{itemize}

\end{document}
